\begin{longtable}{|>{\scriptsize\raggedright\arraybackslash}p{3.5cm}|>{\raggedright\arraybackslash}p{3.5cm}|>{\raggedright\arraybackslash}p{4cm}|>{\raggedright\arraybackslash}p{4cm}|}
\caption{Use Case Test Cases}\label{tab:use_case_tests}\\
\hline
\textbf{Test Case ID} & \textbf{Test Title / Use Case} & \textbf{Test Steps} & \textbf{Expected Results} \\
\hline
\endfirsthead
\multicolumn{4}{c}%
{{\bfseries \tablename\ \thetable{} -- continued from previous page}} \\
\hline
\textbf{Test Case ID} & \textbf{Test Title / Use Case} & \textbf{Test Steps} & \textbf{Expected Results} \\
\hline
\endhead
\hline \multicolumn{4}{|r|}{{Continued on next page}} \\ \hline
\endfoot
\hline
\endlastfoot

% UC-002: Add New Book (Librarian)
\multicolumn{4}{|l|}{\textbf{Sub-Section: UC-002: Add New Book (Librarian) }} \\ \hline
TC\_UC\_BOOK\_001 & Successful addition of a new book (Main Flow) & 
    \begin{itemize}[noitemsep, topsep=0pt, partopsep=0pt, leftmargin=*]
        \item As a logged-in Librarian, navigate to the Book Management page (`/books`).
        \item Click the "New Book" button.
        \item Fill in all required fields (Name, ISBN) with valid data.
        \item Select a valid Author and Genre from the provided lists.
        \item Ensure "Is Available" is set to true.
        \item Submit the form.
    \end{itemize} & 
    \begin{itemize}[noitemsep, topsep=0pt, partopsep=0pt, leftmargin=*]
        \item The system creates a new book record via `POST /api/book/add`.
        \item A success message is displayed (e.g., "Book added").
        \item The "Add Book" form closes.
        \item The new book appears in the list on the Book Management page.
    \end{itemize} \\ \hline
TC\_UC\_BOOK\_002 & Attempt to add a book with missing required fields (Alternative Flow A1) & 
    \begin{itemize}[noitemsep, topsep=0pt, partopsep=0pt, leftmargin=*]
        \item As a logged-in Librarian, navigate to the "Add Book" form.
        \item Fill in the form but leave a required field (e.g., Name) blank.
        \item Submit the form.
    \end{itemize} & 
    \begin{itemize}[noitemsep, topsep=0pt, partopsep=0pt, leftmargin=*]
        \item The system rejects the submission due to a validation error.
        \item An error message is displayed, indicating the required field is missing.
        \item The form remains open for correction.
        \item The book is not added to the database.
    \end{itemize} \\ \hline
TC\_UC\_BOOK\_003 & Attempt to add a book that fails backend validation (Alternative Flow A2) & 
    \begin{itemize}[noitemsep, topsep=0pt, partopsep=0pt, leftmargin=*]
        \item As a logged-in Librarian, navigate to the "Add Book" form.
        \item Fill in all fields with data that is syntactically valid on the client-side but might trigger a server error (e.g., a duplicate ISBN if the server enforces uniqueness).
        \item Submit the form.
    \end{itemize} & 
    \begin{itemize}[noitemsep, topsep=0pt, partopsep=0pt, leftmargin=*]
        \item The system displays a generic error message from the API (e.g., "Something went wrong, please try again").
        \item The book is not added to the database.
    \end{itemize} \\ \hline

% UC-003: Borrow a Book (Member)
\multicolumn{4}{|l|}{\textbf{Sub-Section: UC-003: Borrow a Book (Member) }} \\ \hline
TC\_UC\_BORROW\_001 & Successful borrowing of an available book (Main Flow) & 
    \begin{itemize}[noitemsep, topsep=0pt, partopsep=0pt, leftmargin=*]
        \item As a logged-in Member, navigate to the Borrowal Management page (`/borrowals`).
        \item Click the "New Borrowal" button.
        \item Select a book that is marked as available from the list.
        \item Verify that the Member field is correctly auto-populated.
        \item Submit the form.
    \end{itemize} & 
    \begin{itemize}[noitemsep, topsep=0pt, partopsep=0pt, leftmargin=*]
        \item A new borrowal record is created with "Borrowed" status via `POST /api/borrowal/add`.
        \item The availability of the book is updated to `false`.
        \item A success message is displayed.
        \item The member's borrowal list is updated to show the new borrowal.
    \end{itemize} \\ \hline
TC\_UC\_BORROW\_002 & Attempt to borrow an unavailable book (Alternative Flow A1) & 
    \begin{itemize}[noitemsep, topsep=0pt, partopsep=0pt, leftmargin=*]
        \item As a logged-in Member, navigate to the "Add Borrowal" form.
        \item Attempt to select a book that is not available (isAvailable=false).
    \end{itemize} & 
    \begin{itemize}[noitemsep, topsep=0pt, partopsep=0pt, leftmargin=*]
        \item The UI should ideally prevent the selection of an unavailable book.
        \item If selection is possible, the system should prevent submission or display an error message (e.g., "Book is not available") upon submission attempt.
        \item No borrowal record is created.
    \end{itemize} \\ \hline
TC\_UC\_BORROW\_003 & Verify correct data population in borrowal form & 
    \begin{itemize}[noitemsep, topsep=0pt, partopsep=0pt, leftmargin=*]
        \item As a logged-in Member, open the "Add Borrowal" form.
        \item Observe the date fields.
    \end{itemize} & 
    \begin{itemize}[noitemsep, topsep=0pt, partopsep=0pt, leftmargin=*]
        \item The Member field is correctly auto-populated with the current user's ID.
        \item The Borrowed Date is auto-populated to the current date.
        \item The Due Date is auto-populated to a future date (e.g., 14 days from current).
        \item The initial status is correctly set (e.g., "Borrowed").
    \end{itemize} \\ \hline

% UC-005: View Borrowal History (Member)
\multicolumn{4}{|l|}{\textbf{Sub-Section: UC-005: View Borrowal History (Member) }} \\ \hline
TC\_UC\_HISTORY\_001 & View non-empty borrowal history (Main Flow) & 
    \begin{itemize}[noitemsep, topsep=0pt, partopsep=0pt, leftmargin=*]
        \item Log in as a Member who has previously borrowed one or more books.
        \item Navigate to the Borrowal Management page (`/borrowals`).
    \end{itemize} & 
    \begin{itemize}[noitemsep, topsep=0pt, partopsep=0pt, leftmargin=*]
        \item The system fetches all borrowals and the client filters them.
        \item A list/table is displayed showing *only* the records where the `memberId` matches the logged-in Member's ID.
        \item The displayed data (Book Name, Borrowed Date, Due Date, Status) is accurate for each record.
    \end{itemize} \\ \hline
TC\_UC\_HISTORY\_002 & View empty borrowal history (Alternative Flow A1) & 
    \begin{itemize}[noitemsep, topsep=0pt, partopsep=0pt, leftmargin=*]
        \item Log in as a new Member with no borrowal records.
        \item Navigate to the Borrowal Management page (`/borrowals`).
    \end{itemize} & 
    \begin{itemize}[noitemsep, topsep=0pt, partopsep=0pt, leftmargin=*]
        \item The system displays a message indicating there is no borrowal history (e.g., "No borrowals found").
    \end{itemize} \\ \hline
TC\_UC\_HISTORY\_003 & Verify role-based view of borrowal history & 
    \begin{itemize}[noitemsep, topsep=0pt, partopsep=0pt, leftmargin=*]
        \item Prerequisite: There are borrowals from multiple members in the system.
        \item Log in as a Librarian and navigate to the Borrowal Management page (`/borrowals`).
        \item Log out and log in as a Member (who has at least one borrowal) and navigate to the same page.
    \end{itemize} & 
    \begin{itemize}[noitemsep, topsep=0pt, partopsep=0pt, leftmargin=*]
        \item As the Librarian, all borrowal records for all members are visible.
        \item As the Member, only their own borrowal history is visible.
    \end{itemize} \\ \hline
    
% UC-001 & UC-004 are covered in the Authentication/Authorization test cases document.
\multicolumn{4}{|l|}{\parbox{16cm}{\vspace{5pt}\textbf{Note: UC-001 (User Login) and UC-004 (Register New User) are detailed in the Authentication \& Authorization test plan.}\vspace{5pt}}} 

\end{longtable}
